\section{Scope, notation, and the status of the results}

This article continues the collective manuscript \emph{Polylogarithms and
their Arithmetic Bridges} in Vladimir Reshetnikov's ProveIt repository.
The source review is frozen at commit
\begin{center}\small\ttfamily
570b0567f311cf1890865065896be2665f469e4f.
\end{center}
The two entry points supplied for this task were the canonical manuscript
directory and \texttt{docs/incoming}. At the frozen revision the latter
contains its intake ledger; the relevant previous arrivals have been placed
in the thematic reports tree. The canonical README reports a 379-page
manuscript and explicitly distinguishes preserved reports, finite replay,
and integrated proofs \cite{repoManuscript}. We follow that distinction.

\subsection{What was already available}
The reviewed source includes the established $S_4$ proof, rational
distribution-rank theorems, signed-kernel results, log-Gamma moment
asymptotics, Stieltjes parameter derivatives, and the existing sharp
fixed-argument Herglotz expansion. These form the baseline, rather than new
claims of this continuation. In particular, the exact evaluation of the
manuscript's $S_6$ remains a separate problem. Nothing in the present work
proves its special-value conjecture.

Two previous developments are particularly close to the present analysis.
The boundary continuation already pairs the Gamma and zeta poles in the
Lerch expansion and proves endpoint differentiability thresholds
\cite{repoBoundary}. The integrated zeta-jet chapters already identify
Stieltjes primitives with spectral derivatives of Hurwitz zeta at zero
\cite{repoIntegration}. Our first two analytic developments connect and
extend those mechanisms in different directions: uniformity in a growing
integer order, and nonzero shifts in bilinear integrals.

\subsection{The principal additions}
Table~\ref{scope:contributions} separates the new conclusions proved here from their
classical or previously established inputs. ``New'' in this table means a
continuation beyond the inspected ProveIt material; no claim of worldwide
priority is made without a separate literature investigation.

\begin{table}[htbp]
\centering
\small
\caption{Contributions and their established inputs.}
\label{scope:contributions}
\medskip
\begin{tabular}{@{}>{\raggedright\arraybackslash}p{0.19\textwidth}>{\raggedright\arraybackslash}p{0.44\textwidth}>{\raggedright\arraybackslash}p{0.29\textwidth}@{}}
\toprule
Topic & Conclusion established here & Input credited separately\\
\midrule
Resonant Lerch limits & An arbitrary-order expansion uniform in all integer
$k\geq0$, with a common centered scale and an exponentially small analytic
tail & Classical Lerch expansion; fixed-order pole cancellation\\[5pt]
Shifted Gamma correlations & Exact autocorrelation and increment formulas;
strict convexity of the autocorrelation; unique half-period extremum;
mixed Hurwitz-jet cusps & Hurwitz Fourier expansion; classical unshifted
integral calculus\\[5pt]
Herglotz joint limit & Complete compact-ratio expansion and a finite-$r$
Taylor bound; a profile with arithmetic, Bessel, Gamma, and dilogarithm
representations & Digamma integral; classical partial fractions and zeta
functional equation\\[5pt]
Nonseparable jets & A two-generator colength theorem in any local
Frobenius algebra; exact reflection dimensions and integer torsion;
three-generator obstruction & The manuscript's all-ring distribution
basis and reflection--Koszul model\\
\bottomrule
\end{tabular}
\end{table}

The analytic results have complementary roles. The Lerch theorem controls a
moving spectral order near an integer. The Gamma correlation theorem
integrates two spectral jets against a periodic shift. The Herglotz theorem
controls a moving argument when differentiation order grows. The final
algebraic section concerns the formal distribution relations satisfied by
families of these functions; it makes no assertion of transcendental
independence of their evaluated values.

\subsection{Conventions}
We use the Hurwitz Laurent convention
\begin{equation}\label{scope:laurent}
 \zeta(s,a)=\frac1{s-1}
       +\sum_{n=0}^{\infty}\frac{(-1)^n\gamma_n(a)}{n!}(s-1)^n.
\end{equation}
Thus $\gamma_0(a)=-\psi(a)$ and $\gamma_0(1)=\gamma$, Euler's constant.
Square-bracket superscripts denote spectral derivatives:
\[
 \zeta^{[r]}(s,a)=\partial_s^r\zeta(s,a),\qquad
 \Li_s^{[r]}(z)=\partial_s^r\Li_s(z).
\]
An ordinary prime on $\zeta(s)$ also denotes an $s$ derivative. Parenthesized
superscripts on $\psi$, $\Gamma$, or a named real-variable function denote
derivatives in that function's argument. Each section specifies any further
local notation. All logarithms of positive real numbers are real logarithms;
complex logarithms and square roots are principal unless another branch is
explicitly specified. The fractional part $\{x\}$ is used only under
integrals or away from its integer discontinuities, so values assigned at
the isolated endpoints do not affect the integrals.

The source identities for Hurwitz zeta, Gamma, and Lerch are recorded in
DLMF \cite{dlmfhurwitz,dlmflerch}. In particular,
\begin{equation}\label{scope:lerchidentity}
 \zeta^{[1]}(0,a)=\log\Gamma(a)-\tfrac12\log(2\pi),
 \qquad \partial_a\zeta(s,a)=-s\zeta(s+1,a).
\end{equation}
These identities fix the normalization of every antiderivative below.

\subsection{How the evidence is used}
Theorems in this article are proved analytically or by exact linear algebra.
The computational artifacts have three distinct purposes. Symbolic
recurrences and finite-field ranks are checked exactly. Independent
high-precision evaluations test normalizations and signs in the analytic
identities. Asymptotic tables illustrate the stated remainder orders, without
claiming that a finite table proves a uniform estimate. None of the proofs
depends on accepting a small floating-point residual as zero.
